\section{Inherited HAO Pipeline}

\subsection{Two queries share one CTA}

HAO's central contribution is a complete ownership plan for Blackwell
\cite{hao2026fp4flashattention}. For D128, one 16-warp CTA advances two
independent M128 query tiles, called stage 0 and stage 1. A load warp keeps K
and V ahead in a TMA-fed shared-memory ring. One warp issues ordered QK and PV
operations. Two softmax warpgroups process the query stages in ping-pong
fashion, while separate correction and epilogue roles update online state and
store the result.

Small hardware barriers act as event counters, not block-wide meetings. They
announce ``score ready'', ``first P half ready'', ``tail P ready'', and
``output released''. Algorithm~\ref{alg:hao-pipeline} states the retained
ownership protocol. Operations for $q=0$ and $q=1$ are interleaved whenever
their dependencies allow; the order shown is per score bank, not a serialized
whole-CTA loop.

\begin{algorithm}[htbp]
\caption{Two-query score-to-PV ownership protocol inherited from HAO}
\label{alg:hao-pipeline}
\small
\begin{tabularx}{\textwidth}{@{}r p{.18\textwidth} X@{}}
\toprule
& Owner & Event for query stage $q\in\{0,1\}$ \\
\midrule
1 & Load warp & Prefetch the next K/V tile into the shared-memory ring. \\
2 & MMA warp & Wait until score bank $S_q$ is free; issue QK into $S_q$ and publish \emph{score ready}. \\
3 & Softmax WG $q$ & Read N32 quarters Q0 and Q1, create packed P and scales in retired score addresses, then publish the first legal K64 operand. \\
4 & MMA warp & Observe the first-half event and issue asynchronous $PV_{q,0}$ into permanent output bank $O_q$. \\
5 & Softmax WG $q$ & Process Q2 and Q3 and publish the tail K64 operand. \\
6 & MMA warp & Issue $PV_{q,1}$; meanwhile issue legal work for stage $1-q$. \\
7 & Correction WG & Update online state; the epilogue normalizes and stores when the key loop ends. \\
8 & MMA warp & After tail PV has consumed the overlay, publish $S_q$ as free for the next QK tile. \\
\bottomrule
\end{tabularx}
\end{algorithm}

The two-query topology is essential at high head count: one CTA performs two
query jobs while reusing the same K/V stream. A one-query topology creates
twice as many independent CTAs and can add an extra GPU wave. A third query
stage would provide more look-ahead, but the TMEM layout leaves no bank for it.

\subsection{TMEM is full, so score storage must be recycled}

SM100 and the tested SM103 path expose 512 logical TMEM columns. A D128 FP32
score tile or output accumulator uses 128 columns. Two score banks and two
output banks therefore consume the entire allocation.

\begin{figure}[htbp]
\centering
\begin{tikzpicture}[x=0.044\textwidth,y=0.72cm,font=\small]
  \draw[thick] (0,0) rectangle (20,1);
  \fill[scoreblue!22] (0,0) rectangle (5,1);
  \fill[scoreblue!36] (5,0) rectangle (10,1);
  \fill[outputgreen!26] (10,0) rectangle (15,1);
  \fill[outputgreen!40] (15,0) rectangle (20,1);
  \draw (5,0) -- (5,1);
  \draw (10,0) -- (10,1);
  \draw (15,0) -- (15,1);
  \node at (2.5,.55) {$S_0$: score stage 0};
  \node at (7.5,.55) {$S_1$: score stage 1};
  \node at (12.5,.55) {$O_0$: output stage 0};
  \node at (17.5,.55) {$O_1$: output stage 1};
  \node[below] at (0,0) {0};
  \node[below] at (5,0) {128};
  \node[below] at (10,0) {256};
  \node[below] at (15,0) {384};
  \node[below] at (20,0) {512};
  \node[align=center,text=inkgray] at (5,-1.05)
    {temporary: QK score, then packed $P$ and scale pages};
  \node[align=center,text=inkgray] at (15,-1.05)
    {long lived: FP32 online output accumulators};
\end{tikzpicture}
\caption{D128 TMEM layout inherited from HAO. P and its instruction-facing
scales reuse score storage only after that score fragment is retired. There
is no unowned 128-column bank.}
\label{fig:tmem-layout}
\end{figure}

Each score bank is reused as
\begin{equation}
  \text{QK score}\rightarrow\text{read N32 fragment}
  \rightarrow\text{P/scale overlay}\rightarrow\text{PV consume}
  \rightarrow\text{free}.
  \label{eq:score-overlay}
\end{equation}
Writing the next QK tile too early destroys P still owned by PV. Waiting too
long leaves the matrix issuer idle. The scheduling problem is therefore an
ownership handoff, not merely a shortage of barrier signals. A barrier can
report that storage is free; it cannot create another destination.

\subsection{Quartering creates an early publication point}

QK still writes one M128$\times$N128 FP32 score tile. ``Quartering'' means
that a softmax warpgroup reads and transforms four N32 fragments,
\begin{equation}
  S=[Q0\mid Q1\mid Q2\mid Q3],\qquad
  Qk\in\mathbb R^{128\times32},
\end{equation}
not that QK becomes four smaller matrix products. The score allocation remains
128 columns.

The available scaled-FP4 PV path consumes K64 rather than K32. Two adjacent
quarters must therefore be complete before useful matrix work can issue:
\begin{equation}
  (Q0,Q1)\rightarrow PV_{K64}^{(0)},\qquad
  (Q2,Q3)\rightarrow PV_{K64}^{(1)}.
  \label{eq:k64-pairs}
\end{equation}
Quartering reduces the live register fragment and lets P overwrite retired
score addresses. It does not reduce score-bank TMEM. Its performance value is
the event after Q1: first-half PV can begin while Q2/Q3 and the other query
stage provide independent work.

\begin{figure}[htbp]
\centering
\begin{tikzpicture}[
  x=1.18cm,y=.82cm,font=\small,
  box/.style={draw,rounded corners=1.5pt,minimum width=1.22cm,
              minimum height=.56cm,align=center},
  dep/.style={-{Stealth[length=2mm]},thick}]
  \node[box,fill=scoreblue!26] (s0) at (0,1.5) {QK 0};
  \node[box,fill=probred!17] (a0) at (1.5,1.5) {Q0};
  \node[box,fill=probred!23] (a1) at (2.8,1.5) {Q1};
  \node[box,fill=outputgreen!25,minimum width=1.55cm] (p0) at (4.25,1.5) {PV 0a};
  \node[box,fill=probred!28] (a2) at (5.75,1.5) {Q2};
  \node[box,fill=probred!34] (a3) at (7.05,1.5) {Q3};
  \node[box,fill=outputgreen!38,minimum width=1.55cm] (p1) at (8.5,1.5) {PV 0b};

  \node[box,fill=scoreblue!38] (s1) at (1.0,0) {QK 1};
  \node[box,fill=probred!17] (b0) at (2.5,0) {Q0};
  \node[box,fill=probred!23] (b1) at (3.8,0) {Q1};
  \node[box,fill=outputgreen!25,minimum width=1.55cm] (p2) at (5.25,0) {PV 1a};
  \node[box,fill=probred!28] (b2) at (6.75,0) {Q2};
  \node[box,fill=probred!34] (b3) at (8.05,0) {Q3};
  \node[box,fill=outputgreen!38,minimum width=1.55cm] (p3) at (9.5,0) {PV 1b};

  \draw[dep] (s0)--(a0); \draw[dep] (a0)--(a1);
  \draw[dep] (a1)--(p0); \draw[dep] (a1)--(a2);
  \draw[dep] (a2)--(a3); \draw[dep] (a3)--(p1);
  \draw[dep] (s1)--(b0); \draw[dep] (b0)--(b1);
  \draw[dep] (b1)--(p2); \draw[dep] (b1)--(b2);
  \draw[dep] (b2)--(b3); \draw[dep] (b3)--(p3);
\end{tikzpicture}
\caption{Dependency sketch for the two query stages. Box width is not
measured time. Each first PV half waits for two N32 producer quarters; work
from the other stage covers part of that delay.}
\label{fig:quarter-pipeline}
\end{figure}
